% !TEX root = ../main.tex
%
% Methods subsection -- Datasets and Cohorts
% Requires (in preamble of main.tex): booktabs

\subsection{Datasets and Cohorts}
\label{subsec:cohorts}

Two independent cohorts were used: The Cancer Genome Atlas Pancreatic
Adenocarcinoma cohort (TCGA-PAAD) and the Clinical Proteomic Tumor Analysis
Consortium Pancreatic Ductal Adenocarcinoma cohort (CPTAC-PDAC). A patient was
included if all of the following were available: at least one successfully
processed diagnostic H\&E whole-slide image, overall survival labels (time and
event indicator), age and sex, and an RNA-seq expression profile. Under these
criteria, 152 TCGA-PAAD patients (466 whole-slide images, 2.49 slides per
patient on average) and 144 CPTAC-PDAC patients (564 whole-slide images, 3.31
slides per patient on average) were retained. Patients with more than one slide
contribute all of their slides as a single patient-level bag: slide-level
embeddings are computed independently and mean-pooled into one representation
per patient before the risk head is applied.

Table~\ref{tab:cohort-summary} summarizes the two cohorts. The two institutions
differ not only in scanning/preparation pipeline but also in event composition:
CPTAC-PDAC has a substantially higher observed-death rate (80.6\%) than
TCGA-PAAD (54.6\%), i.e.\ TCGA-PAAD is more heavily censored. This
institution-level heterogeneity, in addition to institution-level batch effects,
motivates evaluating models under true cross-institution external validation
rather than a pooled split of both cohorts.

\begin{table}[htbp]
  \centering
  \begin{tabular}{@{}l c c@{}}
    \toprule
    & \textbf{TCGA-PAAD} & \textbf{CPTAC-PDAC} \\
    \midrule
    Patients ($n$)                      & 152             & 144             \\
    Whole-slide images (total)          & 466              & 564              \\
    Slides per patient (mean)           & 2.49             & 3.31             \\
    Age, years (mean $\pm$ SD)          & 64.8 $\pm$ 10.9  & 64.5 $\pm$ 10.8  \\
    Sex (male / female)                 & 82 / 70          & 75 / 69          \\
    Deaths / censored                   & 83 / 69          & 116 / 28         \\
    Event rate                          & 54.6\%           & 80.6\%           \\
    Median OS, days (IQR)               & 456 (263--658)   & 442 (254--856)   \\
    \bottomrule
  \end{tabular}
  \caption{Cohort summary for the two institutions after applying the inclusion
    criteria (WSI + overall survival label + clinical + RNA-seq all present).
    OS: overall survival; IQR: interquartile range.}
  \label{tab:cohort-summary}
\end{table}
